\uuid{vFVS}
\titre{Équations complexes}
\auteur{bodin ; hanani ; mzari}
\organisation{exo7}
\licence{CC BY-NC-SA 4.0}
\datecreate{2020-01-20}
\niveau{L1}
\module{Algèbre}
\chapitre{Nombres complexes}
\sousChapitre{Racine carrée, équation du second degré}
\difficulte{3}
\theme{Équations, équations complexes}
\exo7qcmid{lille-1.nombres-complexes.026}

\contenu{
\begin{question}
\qclassification{104.02,104.03,104.04}
\qtags{motcle=équations complexes}
On considère l'équation $(E) : \, z^8= \overline{z}, \, z\in \Cc$. Quelles sont les assertions vraies ?
\begin{answers}
    \bad{Si $z$ est une solution de $(E)$, alors $z=0$.}

    \good{Si $z$ est une solution de $(E)$, alors $z=0$ ou $|z|=1$.}

    \bad{L'équation $(E)$ admet $8$ solutions distinctes.}

    \good{Les solutions non nulles de $(E)$ sont les racines $9$-ièmes de l'unité.}

\end{answers}
\begin{explanations}
Remarquer que si $z$ est une solution  de $(E)$, $|z|^8=|\overline{z}|=|z|$, donc si $z$ n'est pas nul,  $|z|=1$.
Par conséquent, $z$ est une solution non nulle de $(E)$ si et seulement si  $z^9=z\overline{z}=1$.
\end{explanations}
\end{question}
}
